% Chapter: Candidate Substrates
% Comparative evaluation of physical computing substrate families.
% NOTE (2026-07-30): the backbone of this chapter (the selection criteria,
% per-family sketches, tab:substrate_comparison, and the verdict) was moved
% out of the Introduction, where it previously formed sec:lit_substrates.
% The label sec:lit_substrates is retained (on the closing section) so that
% existing cross-references keep working.

\section{The Substrate Search}
\label{sec:substrate_search}

% --- EXEMPLAR PARAGRAPH (rewrite in your own words) ---
This project did not begin with a fixed substrate; it began with the Project 11 premise that a fluid or chemical medium should compute continuously and in parallel, and with a search for the medium best suited to that premise. Six candidate families were evaluated: acoustic waves in porous media, free-field acoustic channel geometries, thermal and optothermal media, well-mixed chemical reaction networks, DNA strand-displacement circuits, and reaction--diffusion excitable media. The evaluation criteria were: continuity and parallelism of the computation; trainability (does the medium expose structure that can act as weights?); native nonlinearity; signal speed; input/output practicality; simulation tractability; experimental realisability; and novelty headroom in the literature.
% --- END EXEMPLAR ---

% --- EXEMPLAR PARAGRAPH (rewrite in your own words) ---
This chapter walks through that search family by family. For each candidate it summarises the governing physics, what the literature has achieved with it, what this project itself did with it, and why it was ultimately set aside. \refSec{sec:sub_acoustic} covers the acoustic family, \refSec{sec:sub_thermal} the thermal and optothermal route, \refSec{sec:sub_crn} well-mixed chemistry including DNA strand displacement, and \refSec{sec:sub_rd} the reaction--diffusion excitable media that were finally selected. \refSec{sec:lit_substrates} collects the comparison in \refTab{tab:substrate_comparison} and states the selection rationale.
% --- END EXEMPLAR ---

\section{Acoustic Wave Computing}
\label{sec:sub_acoustic}

% --- EXEMPLAR PARAGRAPH (rewrite in your own words) ---
Sound is an intuitively appealing computing substrate: acoustic waves propagate in ordinary fluids and solids, interact strongly with structured boundaries, and are generated and detected with simple transducers. Two parameterisation styles exist within the family. In free-field and channel geometries, the shape of the domain itself, junctions, labyrinths, metasurfaces, steers interference and scattering. In porous media, the pore architecture modulates the effective density and bulk modulus of the saturating fluid, so that continuous material-property fields (porosity, flow resistivity, tortuosity) offer a trainable design space rather than a purely geometric one.
% --- END EXEMPLAR ---

% --- EXEMPLAR PARAGRAPH (rewrite in your own words) ---
Wave physics is the most mature branch of the physical-neural-network programme reviewed in \refSec{sec:lit_pnn}. Metamaterial slabs have been arranged to perform analog mathematical operations on incident fields \cite{silva2014performing}, and linear wave propagation can be written exactly as an analog recurrent neural network, with the wave field as the hidden state and the local material distribution as the recurrence weights \cite{hughes2019wave}. Deep physical networks can be trained by backpropagation through differentiable simulations of the hardware \cite{wright2022deep}, backpropagation-free in-situ training methods have addressed the noise and differentiability barriers of real wave devices \cite{momeni2023backpropagation}, and a fabricated passive acoustic metamaterial with trained phase-shift elements has performed real-time object recognition \cite{weng2020meta}. Taken together, these results show that wave substrates can be both trained and fabricated, and, by the same token, that the headline demonstrations an MSc project might aim at in this family have largely already been made. A further structural limitation is that passive linear wave systems collapse to a single linear map regardless of depth, so computational nonlinearity must be imported from outside the medium \cite{mcmahon2024nonlinear}.
% --- END EXEMPLAR ---

% --- EXEMPLAR PARAGRAPH (rewrite in your own words) ---
This project investigated the free-field branch directly. Using a finite-difference time-domain solver written by the author, acoustic pulses were routed through a structured Y-junction channel, and their interference at the junction was shown to reproduce elementary logic-gate behaviour at the output probe. The episode established the pulse-in, probe-out experimental grammar, localised inputs, time-series readouts, no inspection of the internal field, that the present reaction--diffusion work retains.
% --- END EXEMPLAR ---

% --- EXEMPLAR PARAGRAPH (rewrite in your own words) ---
The acoustic family was set aside on two of the selection criteria. First, native nonlinearity: at the amplitudes of interest the physics is linear and passive, so the medium shapes signals by superposition only, and every nonlinearity must be imported from the readout or the training procedure. Second, novelty headroom: acoustic and metamaterial logic, routing, and trained classifiers are extensively covered in the literature, and the field's own roadmaps describe it as maturing, leaving little room for an original contribution at MSc scale. The porous variant, although less crowded, couples its most natural trainable knob (flow resistivity) directly to dissipation, the property that destroys the coherent interference the computation depends on.
% --- END EXEMPLAR ---

\section{Thermal and Optothermal Media}
\label{sec:sub_thermal}

% --- EXEMPLAR PARAGRAPH (rewrite in your own words) ---
Heat transport in structured media is diffusive, continuous, and easily interfaced: temperature fields can be written optically and read thermally, and they modulate other physical properties, for example the local sound speed in air, offering an indirect route to programmable wave media. Because the same diffusion equation is cheap to simulate, the thermal route has an attractive simulation-to-effort ratio.
% --- END EXEMPLAR ---

% --- EXEMPLAR PARAGRAPH (rewrite in your own words) ---
The literature records genuine but one-sided progress. Analog matrix computation by topology-optimised heat conduction has been demonstrated in simulation with high accuracy \cite{silva2026thermal}, thermal transistors and thermal logic gates have been proposed within phononics \cite{li2012colloquium}, and thermo-optic waveguide chips have been trained for classification tasks \cite{shen2017deep}. Macroscopic thermal logic, however, has never been realised experimentally despite two decades of theory \cite{li2012colloquium}, and the fundamental scaling is adverse: the thermal time constant grows with the square of device size, so the substrate slows down exactly as it is scaled up.
% --- END EXEMPLAR ---

% --- EXEMPLAR PARAGRAPH (rewrite in your own words) ---
This project tested the thermal route with a candidate thermoviscous mechanism: steady heat patterns were applied along an acoustic channel, with the intention that the temperature-induced modulation of the medium would tune the channel's transfer of acoustic energy. The investigation produced a null result, the relevant input--output frequency ratio proved invariant across all heat patterns applied, removing the very tuning knob the approach depended on.
% --- END EXEMPLAR ---

% --- EXEMPLAR PARAGRAPH (rewrite in your own words) ---
The thermal family was therefore set aside primarily on trainability: the project's own null result removed the design handle, and the external evidence on diffusive slowness and thermal crosstalk corroborates the decision. Weak native nonlinearity and very slow signalling reinforced the verdict; the direction was closed early and not revisited.
% --- END EXEMPLAR ---

\section{Well-Mixed Chemical Reaction Networks}
\label{sec:sub_crn}

% --- EXEMPLAR PARAGRAPH (rewrite in your own words) ---
In a well-mixed chemical computer, the state of the machine is a vector of concentrations evolving in time under mass-action kinetics, and the reaction rates play the role of network weights. Inputs are encoded as initial or injected concentrations, outputs are read from concentration levels or differences, and the ``program'' is the reaction network itself, either designed reaction by reaction, or realised in DNA strand-displacement chemistry where the weights are synthesised directly as nucleotide sequences.
% --- END EXEMPLAR ---

% --- EXEMPLAR PARAGRAPH (rewrite in your own words) ---
This is the most theoretically complete family in the search. Chemical kinetics was shown early to be capable of implementing neural networks and Turing machines, and indeed universal computation \cite{hjelmfelt1991chemical, magnasco1997chemical}; a self-organising chemical reaction network has been shown to act as a reservoir computer, classifying temporal inputs without any designed circuit structure \cite{baltussen2024chemical}; recurrent neural chemical reaction networks have been proved to approximate arbitrary dynamical systems and trained as chemical classifiers \cite{dack2024rncrn}; and DNA-based winner-take-all networks have been scaled experimentally to pattern-recognition tasks of practical size \cite{cherry2018scaling}. Trainability, the weakest point of most physical substrates, is the strongest point of this one.
% --- END EXEMPLAR ---

% --- EXEMPLAR PARAGRAPH (rewrite in your own words) ---
This project engaged with the family on both branches. The author implemented and trained a recurrent neural chemical reaction network of the type proposed by Dack \etal \cite{dack2024rncrn} and verified that it solves the XOR classification task end-to-end, confirming that well-mixed chemistry is genuinely trainable in the neural-network sense. DNA strand-displacement winner-take-all networks were then explored as the more experimentally advanced branch, and set aside: the chemistry is slow and single-use, and the connection to the fluid-mechanical framing of Project 11 is weak.
% --- END EXEMPLAR ---

% --- EXEMPLAR PARAGRAPH (rewrite in your own words) ---
The decisive criterion, however, was continuity and parallelism. A well-mixed reactor has no spatial extent by construction: all information resides in scalar concentrations, computation is neither continuous across space nor parallel across independent signal streams, and there is no geometry to program. The trained-XOR exercise made the limitation concrete, the network learned, but there was nothing spatial to shape, and directly motivated the move to spatially extended media.
% --- END EXEMPLAR ---

\section{Reaction--Diffusion Excitable Media}
\label{sec:sub_rd}

% --- EXEMPLAR PARAGRAPH (rewrite in your own words) ---
Reaction--diffusion excitable media, exemplified by the Belousov--Zhabotinsky reaction reviewed in \refSec{sec:lit_bz}, spread the computation across a continuous two-dimensional space. Supra-threshold perturbations launch self-sustaining chemical pulses that propagate at a characteristic speed, cannot be re-excited behind their refractory tails, and annihilate on collision; walls, gaps, and junctions shape the wave traffic, and in the photosensitive variant a light field $\phi(x,y)$ modulates the excitability itself. The physics, the phenomenology, and the Oregonator and Barkley models are reviewed in \refSec{sec:lit_chemical} and \refSec{sec:lit_bz} and are not repeated here.
% --- END EXEMPLAR ---

% --- EXEMPLAR PARAGRAPH (rewrite in your own words) ---
The literature of this family, reviewed thematically in \refSec{sec:lit_chemical}, delivers exactly the criteria the other families miss: strong native nonlinearity in the excitation threshold, refractoriness, and annihilation of pulses; continuous parallel space in which many wave streams evolve independently; and multiple physically distinct parameterisations, wall geometry, the light-suppression field, anisotropic diffusion, that can serve as weights \cite{toth1995logic, steinbock1996chemical, adamatzky2002xor, sakurai2002design}. Its honest weaknesses are speed (chemical waves travel at millimetres per minute) and the difficulty of cascading gates without electronic transduction.
% --- END EXEMPLAR ---

% --- EXEMPLAR PARAGRAPH (rewrite in your own words) ---
This project has already established the computational footing of the family: the author's two-dimensional reaction--diffusion solver, implementing both the Oregonator and Barkley kinetics, reproduces the canonical excitable-media phenomenology, target waves, wave-front collision and annihilation, gap diffraction, and propagation timing, giving a verified simulation platform on which the characterisation experiments of this thesis are built.
% --- END EXEMPLAR ---

% --- EXEMPLAR PARAGRAPH (rewrite in your own words) ---
The family was selected because it is the only candidate that combines an A-grade native nonlinearity with continuous parallel space, a compact validated simulation model, and a cheap room-temperature experimental path, while still offering genuine novelty headroom: the literature demonstrates hand-designed gates but no systematic characterisation of what a structured substrate natively computes, and no training of its geometry or property fields. The speed weakness is conceded from the outset and treated as characteristic of the substrate rather than disqualifying.
% --- END EXEMPLAR ---

\section{Selection Rationale}
\label{sec:lit_substrates}

% --- EXEMPLAR PARAGRAPH (rewrite in your own words) ---
\refTab{tab:substrate_comparison} collects the six-family comparison against the eight selection criteria. No candidate scores well on everything; each family carries a fatal weakness, and the question is which weakness can be turned into a contribution. Acoustics is fast and tractable but linear and crowded; thermal media are uncrowded but slow, damped, and, in this project's own test, untrainable; well-mixed chemistry is trainable but spaceless; DNA computing is proven but consumable and equally spaceless.
% --- END EXEMPLAR ---

\begin{table}[htbp]
\centering
\caption{Comparison of candidate computing substrate families against the selection criteria. Cell entries are exemplar shorthand; replace with full assessment wording.}
\label{tab:substrate_comparison}
\small
\begin{tabular}{@{}>{\raggedright\arraybackslash}p{2.2cm}>{\raggedright\arraybackslash}p{1.5cm}>{\raggedright\arraybackslash}p{1.5cm}>{\raggedright\arraybackslash}p{1.4cm}>{\raggedright\arraybackslash}p{1.2cm}>{\raggedright\arraybackslash}p{1.4cm}>{\raggedright\arraybackslash}p{1.5cm}>{\raggedright\arraybackslash}p{1.5cm}>{\raggedright\arraybackslash}p{1.3cm}@{}}
\toprule
& \textbf{Continuity / parallelism} & \textbf{Train-ability} & \textbf{Non-linearity} & \textbf{Speed} & \textbf{Input / output} & \textbf{Simulation tractability} & \textbf{Experimental realisability} & \textbf{Novelty headroom} \\
\midrule
Porous acoustic & continuous, parallel & material fields & weak (linear) & fast & transducers & Helmholtz / FDTD & 3D-printed samples & low \\
Free acoustic channels & continuous, parallel & geometry only & weak (linear) & fast & transducers & FDTD (done) & simple rigs & low \\
Thermal / optothermal & continuous, parallel & none found (own null result) & weak & very slow & thermal/optical & diffusion solvers & benchtop & uncertain \\
Well-mixed CRN & scalar only, no space & rates trainable & strong & slow & concentrations & ODEs, easy & hard wet-lab & moderate \\
DNA strand displacement & scalar only, no space & proven (WTA) & strong & very slow & fluorescence & ODEs, easy & demonstrated & low \\
Reaction--diffusion (BZ) & continuous 2D, parallel & geometry, light, anisotropy & strong (excitable) & slow & pulses, imaging & Oregonator (done) & cheap, room temp. & high \\
\bottomrule
\end{tabular}
\end{table}

% --- EXEMPLAR PARAGRAPH (rewrite in your own words) ---
The comparison in \reftab{tab:substrate_comparison} converges on reaction--diffusion excitable media. RD wins on the criteria that matter most for this thesis, native nonlinearity, spatial continuity with parallelism, simulation tractability through a validated compact model, and experimental cost, and it is the only family whose weakest criterion, trainability, is an open research gap rather than a closed verdict: nobody has systematically characterised what a structured RD substrate computes, and nobody has trained its geometry or its light field. Its concession is speed, which is accepted here as a limitation of the substrate rather than a flaw in the premise. The remainder of this thesis therefore treats a structured Oregonator medium as the computing substrate and asks what it natively computes.
% --- END EXEMPLAR ---
